\PassOptionsToPackage{dvipsnames,svgnames,table}{xcolor}
\documentclass[10pt]{article}
\usepackage[a4paper,margin=22mm]{geometry}
\usepackage[no-math]{fontspec}\setmainfont{Latin Modern Roman}
\usepackage{amsmath,amssymb,amsthm,mathtools,graphicx,xcolor,hyperref,xurl,enumitem,mathrsfs}
\setlength{\emergencystretch}{4em}
\newtheorem{Theorem}{Theorem}[section]
\newtheorem{Lemma}[Theorem]{Lemma}
\newtheorem{Corollary}[Theorem]{Corollary}
\newtheorem{Definition}[Theorem]{Definition}
\theoremstyle{definition}
\newtheorem{Notation}[Theorem]{Notation}
\newtheorem{Example}[Theorem]{Example}
\newtheorem{Remark}[Theorem]{Remark}
\theoremstyle{remark}
\newtheorem{Note}[Theorem]{Note}
\DeclarePairedDelimiter\abs{\lvert}{\rvert} 

\DeclareRobustCommand{\magnus}{\genfrac\{\}{0pt}{}}
\def\bk{{\mathbf{k}}}
\def\MM{{\mathsf{M}}}
\def\SS{{\mathsf{S}}}
\def\dep{\mathrm{dep}}
\def\bt{{\mathbf{t}}}
\def\bs{{\mathbf{s}}}
\def\bu{{\mathbf{u}}}
\def\bbr{{\mathbf{r}}}
\def\bkappa{{\boldsymbol \kappa}}
\def\btau{{\boldsymbol \tau}}
\def\brho{{\boldsymbol \rho}}
\def\lp{\langle}
\def\rp{\rangle}
\def\blambda{{\boldsymbol \lambda}}
\def\bmu{{\boldsymbol \mu}}
\def\N{{\mathbb N}}
\def\C{{\mathbb C}}
\def\Z{{\mathbb Z}}
\def\R{{\mathbb R}}
\def\Q{{\mathbb Q}}
\def\bQ{{\overline{\mathbb Q}}}
\def\Gal{{\mathrm{Gal}}}
\def\et{\text{\'et}}
\def\ab{\mathrm{ab}}
\def\proP{{\text{pro-}p}}
\def\padic{{p\mathchar`-\mathrm{adic}}}
\def\la{\langle}
\def\ra{\rangle}
\def\scM{\mathscr{M}}
\def\lala{\la\!\la}
\def\rara{\ra\!\ra}
\def\ttx{{\mathtt{x}}}
\def\tty{{\mathtt{y}}}
\def\ttz{{\mathtt{z}}}
\def\bkappa{{\boldsymbol \kappa}}
\def\scLi{{\mathscr{L}i}}
\def\sLL{{\mathsf{L}}}
\def\cY{{\mathcal{Y}}}
\def\ad{\mathrm{ad\,}}
\def\Spec{\mathrm{Spec}\,}
\def\Ker{\mathrm{Ker}}
\def\CHplus{\underset{\mathsf{CH}}{\oplus}}
\def\check{{\clubsuit}}
\def\kaitobox#1#2#3{\fbox{\rule[#1]{0pt}{#2}\hspace{#3}}\ }
\def\vru{\,\vrule\,}
\newcommand*{\longhookrightarrow}{\ensuremath{\lhook\joinrel\relbar\joinrel\rightarrow}}
\newcommand{\hooklongrightarrow}{\lhook\joinrel\longrightarrow}
\def\nyoroto{{\rightsquigarrow}}
\newcommand{\pathto}[3]{#1\overset{#2}{\dashto} #3}
\newcommand{\pathtoD}[3]{#1\overset{#2}{-\dashto} #3}
\def\dashto{{\,\!\dasharrow\!\,}}
\def\ovec#1{\overrightarrow{#1}}
\def\isom{\,{\overset \sim \to  }\,}
\def\GT{{\widehat{GT}}}
\def\bfeta{{\boldsymbol \eta}}
\def\brho{{\boldsymbol \rho}}
\def\sha{\scalebox{0.6}[0.8]{\rotatebox[origin=c]{-90}{$\exists$}}}
\def\upin{\scalebox{1.0}[1.0]{\rotatebox[origin=c]{90}{$\in$}}}
\def\downin{\scalebox{1.0}[1.0]{\rotatebox[origin=c]{-90}{$\in$}}}
\def\torusA{{\epsfxsize=0.7truecm\epsfbox{torus1.eps}}}
\def\torusB{{\epsfxsize=0.5truecm\epsfbox{torus2.eps}}}

\allowdisplaybreaks
\begin{document}
\section*{1183. Regular-coefficient reconstruction for group-like noncommutative series over a characteristic-zero field}
Hiroaki Nakamura. \emph{Demi-shuffle duals of Magnus polynomials in a free associative algebra}, Algebraic Combinatorics 6(4)(2023), 929--939. \url{https://doi.org/10.5802/alco.287}. Exact official publisher native version; CC BY 4.0, \url{https://creativecommons.org/licenses/by/4.0/}.\par
\paragraph{Selected problem (editorial).} Prove the exact original Theorem 4.1: reconstruct all trailing-X coefficients of a group-like formal series from its regular coefficients and cX. The original characteristic-zero field, constant term one and primitive-letter coproduct conditions remain explicit. Full newly constructed demi-shuffle/Magnus duality, array-binomial inversion, multidepth cancellation and completed tensor/convolution reconstruction remain obligations, not free ordinary prerequisites. Named Magnus elimination, Reutenauer and GKP interfaces remain explicit references. There is one reconstruction outcome; the multi-letter variant and software example are unselected.
\paragraph{Literal source interfaces (editorial).} Native 373 defines the reduced tuple with trailing entry zero; the recursive displays at 377/381 then explicitly multiply its demi-shuffle polynomial by an extra Y. Both literal displays remain beside the exact 265--276 definition and the complete independent author group-by-group coefficient-counting proof at 416--465. No Y deletion or revised recurrence is supplied. Native 777's unbold kappa and 789's parenthesized coefficient shorthand remain with the explicit fixed tuple and full following standard-pairing chain. All general formulas and human proof words are unchanged.
\paragraph{Source formatting (editorial).} Complete native 125--535 and 614--911, including the full alternative cancellation argument, preserve author order, comments and newline guards. The single unselected multi-letter theorem reference uses its exact original printed index 3.2; wrapper counters restore Section 4 and its original equation indices 9--13; the omitted Section 3 contains original equation 8. All mathematics remains editable native TeX. The shuffle glyph is the author's exact native rotated/scaled exists helper, not a substitute or denied package. Original inactive EPS-torus/path helpers are preserved unexpanded in the macro archive; no upstream class or new acquisition. Ordinary standalone typography/pageflow differences are disclosed.
\clearpage
% BEGIN ADAPTED AUTHOR BLOCK125--535
\section{Introduction}

Let $R$ be a commutative integral domain of characteristic 0, and let $R\la X,Y\ra$ be
the free associative algebra generated over $R$ by two (non-commutative) 
letters $X$ and~$Y$.
For $u,v\in R\la X,Y\ra$, we shall write $[u,v]$
to denote the Lie bracket $uv-vu$.
In \cite{M37}, W.~Magnus introduced the associative subalgebra 
$S_X\subset R\la X,Y\ra$ generated by 
(what are called) the elements arising by elimination of $X$:
\begin{equation}
Y^{(0)}:=Y, \quad
Y^{(k+1)}:=[X,Y^{(k)}] \quad
(k=0,1,2,\dots),
\end{equation}
and showed that $S_X$ is freely generated by the $Y^{(k)}$
$(k=0,1,2,\dots)$.
Moreover, he derived that every element $Z$ of  $R\la X,Y\ra$ can be
written uniquely in the form
\begin{equation} \label{1st-elimination}
Z=\alpha_0 X^m +s_1 X^{m-1}+\cdots+s_m,
\end{equation}
where $\alpha_0\in R$, $s_1,\dots, s_m\in S_X$
(see \cite[Hilfssatz 2]{M37}, \cite[Lemma 5.6]{MKS}).
This observation is the first step in the construction 
of the basic Lie elements (an ordered basis of the free Lie algebra), 
which are obtained via repeated elimination, and
whose powered products in decreasing orders give 
the Poincar\'e-Birkoff-Witt
basis of the enveloping algebra~$R\la X,Y\ra$ (\cite[Theorem~5.8]{MKS}).
Apparently, this theory was historically a starting point 
toward subsequent developments of
finer constructions of free Lie algebra bases 
due to Lazard, Hall, Lyndon, Viennot and others
(see, e.g., \cite[Notes 4.5, 5.7]{Reu}). 

In this note, we however stay on the first step of elimination 
\eqref{1st-elimination} and 
look at combinatorial properties of a certain basis 
$\{\MM^{(\bk)}\}_{\bk\in \N_0^{(\infty)}}$ of $R\la X,Y\ra$
(to be called the Magnus polynomials below)
designed as follows: 

\begin{Notation} 
Let $\N_0$ denote the set of non-negative integers, and 
let 
$$
\N_0^{(\infty)}:=\bigcup_{d=0}^\infty
\left(
\prod_{k=1}^{d}\N_0\right)
\times \N_0 
$$ 
be the collection of finite sequences 
$
\bk=(k_1,\dots,k_d; k_\infty)
$
of non-negative integers equipped 
with a special last entry $k_\infty\in \N_0$. 
Here, we consider $(;k_\infty)$ also as elements of 
$\N_0^{(\infty)}$ coming from $d=0$.
For $\bk\in \N_0^{(\infty)}$, define
$|\bk|:=\sum_{i=1}^\infty k_i=k_1+\cdots+k_d+k_\infty$ 
(resp. $\dep(\bk):=d$),
and call it the size (resp. depth) of $\bk$.
\end{Notation}

\begin{Definition}[Magnus polynomial]
\label{defMagnus}
For $\bk=(k_1,\dots,k_d; k_\infty)\in\N_0^{(\infty)}$, define
$$
\MM^{(\bk)}:=Y^{(k_1)}\cdots Y^{(k_d)} \cdot X^{k_\infty} \in R\la X,Y\ra.
$$
We also set $\MM^{(;0)}=1$, and $\MM^{(;k)}=X^k$ for $k=1,2,\dots$.
Note that $\MM^{(k;0)}=Y^{(k)}$ for~$k\ge 0$.
\end{Definition}

\begin{Example} 
$\MM^{(1,0;2)}=Y^{(1)}Y^{(0)}X^2=(XY-YX)YX^2=XY^2X^2-YXYX^2$.
\end{Example}

It is not difficult to see that the Magnus polynomial $\MM^{(\bk)}\in R\la X,Y\ra$ 
is homogeneous
of bidegree $(|\bk|,\dep(\bk))$ in $X$ and $Y$.

The above mentioned Magnus expression \eqref{1st-elimination}
can then be rephrased as
\begin{equation} \label{Magnus}
Z=\sum_{\bk\in \N_0^{(\infty)}}
\alpha_{\bk}\, \MM^{(\bk)}
\end{equation}
with uniquely determined coefficients $\alpha_{\bk}\in R$ for any given 
$Z\in R\la X,Y\ra$.
In other words, the collection $\{\MM^{(\bk)}\mid \bk\in  \N_0^{(\infty)}\}$ 
forms an $R$-linear basis of $R\la X,Y\ra$.

Below in \S 2, we will construct another $R$-linear basis 
$\{\SS^{(\bk)}\mid \bk\in  \N_0^{(\infty)}\}$
(formed by what we call the ``demi-shuffle'' polynomials)
and show that 
$\{\MM^{(\bk)}\}_\bk$
and $\{\SS^{(\bk)}\}_\bk$
are dual to each other under the standard pairing
with respect to the monomials of~$R\la X,Y\ra$
(Theorem \ref{orthogonality}).
We then in \S 3 shortly generalize the duality to the case of
free associative algebras of more variables (Theorem 
\textup{3.2}). 
In \S 4, we apply the formation of dual basis 
to derive a formula of Le--Murakami, Furusho type 
that expresses arbitrary coefficients of a group-like series 
$J\in R\lala X,Y\rara$ in terms of 
the ``regular'' coefficients of $J$
(Theorem \ref{LeMuFu}).  

\section{Demi-shuffle duals and array binomial coefficients}

Let $W$ be the subset of $R\la X,Y\ra$ formed by the monomials in $X,Y$ 
together with $1$, and call any element of $W$ a word. 
It is clear that $W$ forms a free monoid by the concatenation product
that restricts the multiplication of $R\la X,Y\ra$.
Each element of $R\la X,Y\ra$ is an $R$-linear combination of words in $W$.
For two elements $u,v\in R\la X,Y\ra$, define the standard pairing $\lp u,v\rp\in R$
so as to extend $R$-linearly the Kronecker symbol $\lp w,w'\rp:=\delta_{w}^{w'}\in\{0,1\}$
for words $w,w'\in W$.

\begin{Notation} We use the notation $w_\bk:=X^{k_1}Y\cdots X^{k_d}YX^{k_\infty}$
and call it the word associated to $\bk=(k_1,\dots,k_d; k_\infty)\in\N_0^{(\infty)}$.
The mapping $\bk\mapsto w_\bk$ gives a bijection from $\N_0^{(\infty)}$ onto $W$.
(Note that $w_{(;0)}=1$.)
The standard pairing $\lp w_\bk,w_{\bk'}\rp$ is equal to~$0$ or $1$ 
according to whether $\bk\ne \bk'$ or $\bk=\bk'$.
\end{Notation}


The purpose of this section is to describe the dual of the 
Magnus basis
$\{\MM^{(\bk)}\}_{\bk\in \N_0^{(\infty)}}$ with respect to the standard
pairing.

 
\begin{Definition}[Demi-shuffle polynomial]
\label{demi-shuffle}
For $\bk=(k_1,\dots,k_d; k_\infty)\in\N_0^{(\infty)}$, define
$$
\SS^{(\bk)}:=
(\cdots((X^{k_1}Y)\sha X^{k_2})Y)\sha\cdots)\sha  X^{k_d})Y)\sha
X^{k_\infty} \in R\la X,Y\ra,
$$
where $\sha$ denotes the usual shuffle product. 
We also set $\SS^{(;0)}=1$, and $\SS^{(;k)}=X^k$ for $k=1,2,\dots$.
Note that $\SS^{(k;0)}=X^kY$ for $k\ge 0$.
\end{Definition}

The construction of $\SS^{(\bk)}$ can be interpreted 
as forming the linear sum
of all words obtained from 
the word
$w_\bk=X^{k_1}Y\cdots X^{k_d}YX^{k_\infty}$
by consecutively applying ``left shuffles'' of the $X$ letters
and ``concatenations'' of the $Y$ letters in $w_\bk$.

\begin{Example}
Here are a few examples:
$\SS^{(0,1;0)}=(Y\sha X)Y=YXY+XYY$;
$\SS^{(1,1;0)}=((XY)\sha X)Y=XYXY+2XXYY$;
$\SS^{(1,0,1;0)}=(((XY)Y)\sha X)Y=XYYXY+XYXY^2+2X^2Y^3$.
Using the first identity, one can also compute
\begin{align*}
&\SS^{(0,1;1)}=((Y\sha X)Y)\sha X=(YXY+XYY)\sha X \\
&=(YXYX+2YXXY+XYXY)+(XYYX+XYXY+2XXYY) \\
&=2XXYY+2XYXY+XYYX+2YXXY+YXYX.
\end{align*}
\end{Example}

\begin{Theorem}[Duality] \label{orthogonality}
For $\bt, \bk\in \N_0^{(\infty)}$, we have
$$
\lp \SS^{(\bt)},\MM^{(\bk)}\rp=\delta_\bt^{\bk}.
$$
Here $\delta_\bt^{\bk}$ is the Kronecker symbol, i.e., designating $0$ or $1$
according to whether $\bt\ne\bk$ or $\bt=\bk$ respectively.
\end{Theorem}

Before going to the proof of the above theorem, we introduce the following
notation.

\begin{Definition}[Array binomial coefficient]
For $\bt,\bk\in  \N_0^{(\infty)}$ with $\dep(\bt)=\dep(\bk)$, 
$|\bt|=|\bk|$, define
\begin{equation} \label{ArrayBino}
\binom{\bt}{\bk}:=
\binom{t_1}{k_1}
\binom{t_1+t_2-k_1}{k_2}
\cdots
\binom{t_1+\cdots t_d-k_1-\cdots-k_{d-1}}{k_d},
\end{equation}
where $\bt=(t_1,\dots,t_d,t_\infty)$,
$\bk=(k_1,\dots,k_d,k_\infty)$.
We understand $\binom{\bt}{\bk}=1$ if
$\bt=\bk=(;N)$ for some $N\in \N_0$.
We set $\binom{\bt}{\bk}:=0$ if either $\dep(\bt)\ne\dep(\bk)$ or 
$|\bt|\ne|\bk|$ holds.
\end{Definition}

\begin{Remark} 
\label{rem2.6}
The special case $\binom{N,0,\dots,0;0}{k_1,k_2,\dots,k_d;k_\infty}$ 
is the same as the usual multinomial coefficient 
$\binom{N}{k_1,k_2,\dots,k_d,k_\infty}$ in combinatorics.
Note also that $\binom{\bt}{\bk}\ne 0$ implies $t_\infty\le k_\infty$,
as the last factor of $\binom{\bt}{\bk}$ could survive only when 
$(t_1+\cdots t_d-k_1-\cdots-k_{d-1})-k_d=k_\infty-t_\infty\ge 0$.
\end{Remark}

It turns out that the array binomial coefficients give
the expansion of $\SS^{(\bt)}$ as a linear sum of the monomials in $W$.
Recall that, for $\bt=(t_1,\dots,t_d,t_\infty)\in  \N_0^{(\infty)}$, 
$w_\bt$ denotes the word $X^{t_1}YX^{t_2}Y\cdots X^{t_d}YX^{t_\infty}\in W$.



\begin{Lemma}[Monomial expansion]
\label{dS-mono}
$$
\SS^{(\bk)}=\sum_{\bt \in\N_0^{(\infty)}}
\binom{\bt}{\bk} w_{\bt}.
$$
\end{Lemma}
\begin{proof}
Without loss of generality, it suffices to show 
$\la w_{\bt}, \SS^{(\bk)} \ra=\binom{\bt}{\bk}$
in the case $(N:=)\,|\bt|=|\bk|$ and $(d:=)\,\dep(\bt)=\dep(\bk)$.
The assertion is trivial when $d=0$, as then
$\bk=\bt=(;N)$, $\SS^{(\bk)}=X^N=w_\bt$ and $\binom{\bt}{\bk}=1$.
For $d>0$, we argue by induction on $d$.
Suppose $d=1$, $\bk=(k_1;k_\infty)$ and $\bt=(t_1;t_\infty)$. Then
$$
\SS^{(\bk)}=
(X^{k_1}Y)\sha X^{k_\infty}
=\sum_{i=0}^{k_\infty}(X^{k_1}\sha X^i)YX^{k_\infty-i}
=\sum_{i=0}^{k_\infty}\binom{k_1+i}{k_1}
X^{k_1+i}YX^{k_\infty-i}.
$$
Since $N=k_1+k_\infty=t_1+t_\infty$, we have
$\la  w_{\bt}, \SS^{(\bk)} \ra=\binom{k_1+k_\infty-t_\infty}{k_1}
=\binom{t_1}{k_1}$.
Suppose $d>1$ with $\bk=(k_1,\dots,k_d;k_\infty)$ and
$\bt=(t_1,\dots,t_d;t_\infty)$.
Write $\bk'=(k_1,\dots,k_{d-1};0)\in \N_0^{(\infty)}
$. Then
\begin{align*}
\SS^{(\bk)}
&=(((\SS^{(\bk')}Y)\sha X^{k_d})Y)\sha X^{k_\infty} 
\\
&=\sum_{i=0}^{k_\infty}
\SS^{(\bk')}
Y(X^{k_d}\sha X^i)YX^{k_\infty-i} 
\qquad (\text{associativity of $\sha$})
\\
&=\sum_{i=0}^{k_\infty}
\sum_{\bt'}\binom{\bt'}{\bk'}
\binom{k_d+i}{k_d}
w_{\bt'}
X^{k_d+i}YX^{k_\infty-i},
\end{align*}
where $\bt'=(t_1',\dots,t_{d-1}';t_d')\in \N_0^{(\infty)}$ runs over
those tuples 
with $t_1'+\cdots+t_d'=|\bk'|$ %k_1+\cdots+k_{d-1}$
so that $\SS^{(\bk')}$ is expressed as
$\sum_{\bt'}\binom{\bt'}{\bk'}w_{\bt'}$
by the induction hypothesis on $\dep(\bk')=d-1$.
The coefficient of $w_\bt$ in $\SS^{(\bk)}$
can be found in the above summand where
$k_\infty-i=t_\infty$,
$t_d'+k_d+i=t_d$ and
$t_s'=t_s$ ($s=1,\dots,d-1$), hence
$$
\la w_\bt, \SS^{(\bk)} \ra=
\binom{t_1}{k_1}
\cdots
\binom{t_1+\cdots t_{d-1}-k_1-\cdots-k_{d-2}}{k_{d-1}}
\cdot
\binom{k_d+k_\infty-t_\infty}{k_d}.
$$
Since $N=|\bk|=|\bt|$, we have
$k_d+k_\infty-t_\infty
=t_1+\cdots t_{d}-k_1-\cdots-k_{d-1}$.
This establishes the formula  
$\la w_\bt, \SS^{(\bk)} \ra=\binom{\bt}{\bk}$.
\end{proof}

\begin{Remark}
It would be worth noting that 
Lemma \ref{dS-mono} 
can be derived from
counting~$\la w_{\bt}, \SS^{(\bk)} \ra$ 
as the number of certain shufflings of letters in
$w_\bk=X^{k_1}Y\cdots X^{k_d}YX^{k_\infty}$
to produce 
$w_\bt=X^{t_1}Y\cdots X^{t_d}YX^{t_\infty}$.
Assume 
$|\bt|=|\bk|$ and $\dep(\bt)=\dep(\bk)$,
and consider letters $Y$ as partitions between 
groups of letters $X$ in $w_\bk$ and in $w_\bt$. 
Then~$\la w_{\bt}, \SS^{(\bk)} \ra$ is the number of ways
of moving some letters $X$ in $w_\bk$
to the left (beyond any number of $Y$'s)
to form the word $w_\bt$
without changing orders between~$X$'s from the same group
in $w_\bk$.
We count this number 
by enumerating branches of possibilities
for choosing
places of~$X$'s in $w_\bt$ for those moved from 
$w_\bk$ group by group. 
The first binomial factor $\binom{t_1}{k_1}$
of \eqref{ArrayBino} is the number of ways to choose~$k_1$ places for~$X$'s (coming from 
the first group in $w_\bk$) in the first group $X^{t_1}Y$ of $w_\bt$.
The second binomial factor $\binom{t_1+t_2-k_1}{k_2}$ of \eqref{ArrayBino} 
represents the number of ways to choose~$k_2$ 
places for~$X$'s (coming from the second group $YX^{k_2}Y$ in~$w_\bk$)
 in the first two groups $X^{t_1}YX^{t_2}Y$ of~$w_\bt$ 
where the already occupied $k_1$ places in the previous step
are prohibited from being chosen.
We continue the process in the same way. For each given 
$i\in\{2,\dots,d\}$, suppose that destinations of $X$'s 
in~$X^{t_1}Y\cdots X^{t_{i-1}}$ 
from $X^{k_1}Y\cdots YX^{k_{i-1}}Y$ have already been
chosen. Then, the $i$-th binomial
factor $\binom{t_1+\cdots t_i-k_1-\cdots-k_{i-1}}{k_i}$ 
of \eqref{ArrayBino}
represents the number of ways to choose $k_i$ places for 
$X$'s (coming from the $i$-th group $YX^{k_i}Y$ in~$w_\bk$)
in~$X^{t_1}Y\cdots YX^{t_i}Y$ (the first $i$ groups of~$w_\bt$):
there are $t_1+\cdots+t_i$ places for $X$ in~$X^{t_1}Y\cdots YX^{t_i}Y$,
but already $k_1+\cdots+k_{i-1}$ places are occupied 
by earlier choices.
Performing the process until $i=d$ 
verifies the desired identity
$\la w_\bt, \SS^{(\bk)} \ra=\binom{\bt}{\bk}$.
\end{Remark}


\begin{proof}[Proof of Theorem \ref{orthogonality}]
From the formula~$Y^{(k)}=\sum_{i=0}^k (-1)^i\binom{k}{i}X^{k-i}YX^i$~(\cite[(4)]{M37}), 
it is not difficult to see that the expansion of the Magnus polynomial in monomials
is given by 
\begin{equation} \label{Mag-Mono}
\MM^{(\bk)}=\sum_{\bt \in\N_0^{(\infty)}}
\magnus{\bk}{\bt} w_{\bt} 
\end{equation}
with
\begin{equation}
\label{MagnusBino}
\magnus{\bk}{\bt}:=(-1)^{\sum_{i=1}^d (d-i+1)(k_i-t_i)}
\binom{k_1}{k_1-t_1} 
\binom{k_2}{k_1+k_2-t_1-t_2} \cdots
\binom{k_d}{\sum_{i=1}^d(k_i-t_i)}
\end{equation}
for $\bt:=(t_1,\dots, t_d;t_\infty)$ and $\bk:=(k_1,\dots,k_d;k_\infty)$.
Since we have $\lp\SS^{(\bt)},\MM^{(\bk)}\rp=\sum_{\bu\in\N_0^{(\infty)}}
\lp\SS^{(\bt)},w_\bu\rp \lp\MM^{(\bk)},w_\bu\rp$, it suffices to show
\begin{equation} \label{target}
\sum_\bu \magnus{\bk}{\bu} \binom{\bu}{\bt}=\delta_\bk^\bt.
\end{equation}
Noting that a non-zero pairing $\lp\SS^{(\bt)},\MM^{(\bk)}\rp$ occurs only when
$|\bt|=|\bk|$ and $\dep(\bt)=\dep(\bk)$, we may assume without loss of generality 
that $\bu$ in the above summation also runs over those
with the fixed size $N:=|\bt|=|\bk|$ and depth $d:=\dep(\bt)=\dep(\bk)$.
Then, the summation $\sum_\bu$ with $\bu=(u_1,\dots,u_d;u_\infty)$ 
has $d$ independent parameters $u_1,\dots, u_d$ that determine 
$u_\infty=N-\sum_{i=1}^d u_i$. We may also regard each $u_i$ as running over $\Z$,
as the coefficients $\magnus{\bk}{\bu}$, $\binom{\bu}{\bt}$ vanish when 
their combinatorial meaning is lost.
Then, in the summation $\sum_{(u_1,\dots,u_d)\in\Z^d}$ in \eqref{target},
the partial factor of summation 
involved with the last parameter $u_d$ can be factored out in the form:
\begin{align*}
&
\sum_{u_d\in\Z}(-1)^{-u_d}
\binom{k_d}{u_d+\sum_{i=1}^{d-1}(u_i-k_i)}
\binom{u_d+\sum_{i=1}^{d-1}(u_i-t_i)}{t_d} \\
&=
(-1)^{\sum_{i=1}^{d-1}(u_i-k_i)-k_d}
\binom{\sum_{i=1}^{d-1}(k_i-t_i)}{t_d-k_d}.
\end{align*}
(Use \cite[(5.24)]{GKP}.)
Repeating this process inductively on $d$, we eventually find
$$
\lp\SS^{(\bt)},\MM^{(\bk)}\rp=
\binom{0}{t_1-k_1}
\binom{k_1-t_1}{t_2-k_2}
\binom{k_1+k_2-t_1-t_2}{t_3-k_3}
\cdots
\binom{\sum_{i=1}^{d-1}(k_i-t_i)}{t_d-k_d}
$$
which is equal to $\delta^\bk_\bt$ as desired.
\end{proof}


\begin{Corollary} \label{coro}
Each element $u\in R\la X,Y\ra$ can be written as 
$$
u=\sum_{\bk \in\N_0^{(\infty)}}
\lp\SS^{(\bk)}, u\rp \,\MM^{(\bk)}
=\sum_{\bk \in\N_0^{(\infty)}}\lp\MM^{(\bk)}, u\rp \,\SS^{(\bk)}.
$$
\end{Corollary}

Note that only a finite number of summands are nonzero
in either summation above. 
% END AUTHOR BLOCK125--535
\setcounter{section}{3}
\setcounter{equation}{8}
% BEGIN EXACT AUTHOR BLOCK614--911
\section{Application to a formula of Le--Murakami and Furusho type}

In this section, we assume that $R$ is a field and consider 
$R\la X,Y\ra$ as a subalgebra of the ring of non-commutative
formal power series $R\lala X,Y\rara$, where a standard
comultiplication $\Delta$ is defined by setting
$\Delta(a)=1\otimes a+a\otimes 1$ for $a\in\{X,Y\}$.  
An element $J\in R\lala X,Y\rara$ is called group-like if
it has constant term 1 and satisfies $\Delta(J)=J\otimes J$. 
There are many group-like elements: 
for example, the subgroup multiplicatively generated by $\exp(X)$ and $\exp(Y)$
in $R\lala X,Y\rara^\times$ consists of group-like elements and forms a free group of rank 2. 

\begin{Theorem}[Le--Murakami, Furusho type formula]
\label{LeMuFu}
Let $J\in R\lala X,Y\rara$ be a group-like element in the form
$$
J=\sum_{\bk\in\N_0^{(\infty)}} c_\bk w_\bk,
$$
and write $c_X$ for the coefficient $c_{(;1)}$ of $X$ in $J$. 
Then,
$$
c_{(k_1,\dots,k_d;k_\infty)}= \hspace{-0.2cm} \sum_{\substack{s,t\ge 0 \\ s+t=k_\infty}}
(-1)^s
\frac{{(c_X)}^t}{t!}
\hspace{-0.2cm} \sum_{\substack{ s_1,\dots,s_d\ge 0 \\ s=s_1+\dots+s_d}}
\binom{k_1+s_1}{k_1}\cdots \binom{k_d+s_d}{k_d}
c_{(k_1+s_1,\dots,k_d+s_d;0)}.
$$
\end{Theorem}


We first prove an elementary identity that will be used for the proof of the above formula.

\begin{Lemma} \label{lem2}
Let $\bkappa=(k_1,\dots,k_d)\in\N_0^d$ and $\bs=(s_1,\dots,s_d)\in \Z^d$
satisfy $s=s_1+\cdots+s_d\ge 0$ and $k_i+s_i\ge 0$ $(i=1,\dots,d)$.
Then,  we have
$$
\sum_{\btau\in\N_0^d} 
\lp\SS^{(\btau;0)},w_{(\bkappa+\bs;0)}\rp
\cdot \lp\MM^{(\btau;0)}, w_{(\bkappa;s)}\rp
=(-1)^s
\binom{k_1+s_1}{k_1}\cdots \binom{k_d+s_d}{k_d}.
$$
\end{Lemma}

\begin{proof}
We shall compute the left-hand side explicitly
as the sum over 
$\btau\in\N_0^d$ satisfying
$\sum_{i=1}^d t_i=\sum_{i=1}^d (k_i+s_i)$
with
$$
\la\SS^{(\btau;0)},w_{(\bkappa+\bs;0)}\ra=\binom{(\bkappa+\bs;0)}{(\btau;0)}
=\binom{k_1+s_1}{t_1} 
\cdots
\binom{\sum_{i=1}^{d-1}(k_i+s_i)-\sum_{i=1}^{d-2}t_i}{t_{d-1}}
\binom{t_d}{t_d}
$$
by Lemma \ref{dS-mono} and with
\small $$
\la\MM^{(\btau;0)}, w_{(\bkappa;s)}\ra
=\magnus{(\btau;0)}{(\bkappa;s)}
=(-1)^{s+\sum_{i=1}^{d-1}(d-i)(t_i-k_i)}
\binom{t_1}{t_1-k_1}
\cdots
\binom{t_{d-1}}{\sum_{i=1}^{d-1}(t_i-k_i)}
\binom{t_d}{s}
$$ \normalsize
by \eqref{Mag-Mono} and
$s=\sum_{i=1}^{d}(t_i-k_i)$. 
Note that,
since $\binom{(\bkappa+\bs;0)}{(\btau;0)}\magnus{(\btau;0)}{(\bkappa;s)}\ne 0$ 
only when all entries of $\btau=(t_1,\dots,t_d)$ are nonnegative
and $t_1+\cdots+t_d=\sum_{i=1}^d (k_i+s_i)$ (a constant),
the above sum
can be taken over the tuples $(t_1,\dots,t_{d-1})\in\Z^{d-1}$
with entries running as independent integers.
Then, the partial summation involved with the last variable 
$t_{d-1}$ may be factored out as
\small \begin{align*}
&\sum_{t_{d-1}}(-1)^{t_{d-1}}
\binom{\sum_{i=1}^{d-1}(k_i+s_i)-\sum_{i=1}^{d-2}t_i}{t_{d-1}}
\binom{t_{d-1}}{\sum_{i=1}^{d-1}(t_i-k_i)}
\binom{t_d}{s} \\
&=\sum_{t_{d-1}}(-1)^{t_{d-1}}
\binom{\sum_{i=1}^{d-1}(k_i+s_i)-\sum_{i=1}^{d-2}t_i}{ \sum_{i=1}^{d-1}k_i-\sum_{i=1}^{d-2}t_i }
\binom{\sum_{i=1}^{d-1} s_i}{\sum_{i=1}^{d-1}(t_i-k_i)}
\binom{\sum_{i=1}^d(k_i+s_i)-\sum_{i=1}^{d-1}t_i}{\sum_{i=1}^{d} s_i} \\
&=\binom{\sum_{i=1}^{d-1}(k_i+s_i)-\sum_{i=1}^{d-2}t_i}{\sum_{i=1}^{d-1} s_i}
(-1)^{\sum_{i=1}^{d-1}k_i -\sum_{i=1}^{d-2} t_i}
\binom{k_d+s_d}{s_d},
\end{align*} \normalsize
where \cite[(5.21)]{GKP} is applied for the first equality and 
\cite[(5.24)]{GKP} for the second.
After factoring out the constant $\binom{k_d+s_d}{s_d}$ and repeating
the similar process with the other variables $t_{d-2},\dots, t_1$ consecutively,
we eventually obtain the asserted formula.
Below in Note \ref{Quest4.3}, we also provide an alternative proof of the lemma 
free from intricate use of 
\cite[(5.21),(5.24)]{GKP}.
\end{proof}

\begin{proof}[Proof of Theorem \ref{LeMuFu}]
We argue in the beautiful framework exploited in Reutenauer's book \cite[1.5]{Reu}
using the complete tensor product
$$
\mathscr{A}=R\lala X,Y\rara \bar\otimes R\lala X,Y\rara
$$
equipped with a product induced from 
the shuffle product (resp. the concatenation product)
on the left (resp. right) of $\bar\otimes$. 
Recall that the ring of
$R$-linear endomorphisms $\mathrm{End}_RR\lala X,Y\rara$
can be embedded into $\mathscr{A}$ by 
$f\mapsto \sum_{w\in W} w\otimes f(w)$, and that the product of 
$\mathscr{A}$ restricts to the convolution product of $\mathrm{End}_RR\lala X,Y\rara$
defined by $f\ast g:=\mathrm{conc}\circ(f\otimes g)\circ \Delta$
(where ``$\mathrm{conc}$'' means concatenation of left and right sides of $\otimes$).
Note that, for $f\in\mathrm{End}_R R\lala X,Y\rara$ and  $J\in R\lala X,Y\rara$,
we have $f(J)=\sum_{w\in W} \lp w,J\rp f(w)$.

Since, by Corollary \ref{coro}, every word $w$ can be written as 
$\sum_{\bt\in \N_0^{(\infty)}}
\lp \SS^{(\bt)},w \rp \MM^{(\bt)}$, 
the element of $\mathscr{A}$ corresponding to 
the identity $\mathrm{id}\in \mathrm{End}_RR\lala X,Y\rara$ is:
\begin{align*}
\sum_{w\in W} w\otimes w &=
\sum_w w\otimes \sum_\bt \lp \SS^{(\bt)},w \rp\MM^{(\bt)} 
=
\sum_\bt (\sum_w  \lp\SS^{(\bt)},w\rp w) \otimes \MM^{(\bt)} \\
&= 
\sum_\bt \SS^{(\bt)} \otimes \MM^{(\bt)} \\
&=
\left(\sum_{d=0}^\infty 
\sum_{\btau\in\N_0^d} \SS^{(\btau;0)}\otimes \MM^{(\btau;0)}\right)
\cdot
\left(\sum_{t=0}^\infty X^t\otimes X^t\right),
\end{align*}
where we have used
$\SS^{(\bt)}=\SS^{(\btau;t)}=\SS^{(\btau;0)}\sha X^t$ and 
$\MM^{(\bt)}=\MM^{(\btau;t)}=\MM^{(\btau;0)}\cdot X^t$.
Observing that both factors of the last expression above correspond to specific 
$R$-linear endomorphisms,
we can apply $\mathrm{id}$ to $J$ as the convolution product of them and find from
$\Delta(J)=J\otimes J$ that
\begin{equation} \label{J-expansion}
J=\mathrm{id}(J)=
\left(\sum_{d=0}^\infty 
\sum_{\btau\in\N_0^d} \lp \SS^{(\btau;0)},J \rp \MM^{(\btau;0)}\right)
\left(\sum_{t=0}^\infty \frac{(c_X)^t}{t !} X^t\right).
\end{equation}
Note here that the pairing of $J$ with 
$X^t=X^{\sha\, t}/t!$ is equal to $(c_X)^t/t !$,
as easily seen from the fact that the specialization 
$J(X,0)\in R\lala X\rara$ at $Y=0$ is a group like
element $\exp(c_X\cdot X)$.
To complete the proof of Theorem \ref{LeMuFu},
given a fixed $\bk=(\bkappa;k_\infty)=(k_1,\dots,k_d; k_\infty)\in\N_0^{(\infty)}$
and $0\le s\le k_\infty$, 
we compute
the coefficient of~$w_{(\kappa;s)}=X^{k_1}Y\cdots X^{k_d}YX^s$ in the expansion 
of the first factor of the above right-hand side as follows:
\begin{align*}
&\sum_{d=0}^\infty 
\sum_{\btau\in\N_0^d} 
\lp \SS^{(\btau;0)},J \rp 
\lp \MM^{(\btau;0)},w_{(\bkappa;s)} \rp 
=\left\lp\sum_{d=0}^\infty 
\sum_{\btau\in\N_0^d} \lp\SS^{(\btau;0)},J \rp \MM^{(\btau;0)},
w_{(\bkappa;s)} \right\rp \\
=&\left\lp\sum_{d=0}^\infty 
\sum_{\btau\in\N_0^d} 
\biggl\lp\SS^{(\btau;0)},\sum_{\bu\in \N_0^{(\infty)}}(J,w_\bu)w_\bu\biggr\rp
 \MM^{(\btau;0)},
w_{(\bkappa;s)}\right\rp \\
=&\sum_{\bu} \lp J,w_\bu \rp \sum_{d=0}^\infty 
\sum_{\btau\in\N_0^d} \lp \SS^{(\btau;0)},w_\bu \rp
\lp \MM^{(\btau;0)}, w_{(\bkappa;s)} \rp .
\end{align*}
But since $\lp\SS^{(\btau;0)},w_\bu\rp \lp \MM^{(\btau;0)}, w_{(\bkappa;s)} \rp$ 
survives only when 
$\dep(\btau;0)=\dep(\bkappa;s)=\dep(\bu)$ and $|(\btau;0)|=|(\bkappa;s)|=|\bu|$,
the summation $\sum_\bu$ in the last expression above occurs only for 
those $\bu$ of the form 
$(\bkappa+\bs;0)\in \N_0^{(\infty)}$ 
with $\bs=(s_1,\dots,s_d)\in\Z^d$, 
$s=s_1+\cdots+s_d\ge 0$
(cf. also Remark \ref{rem2.6}).  
Then, it follows from Lemma \ref{lem2} that the last expression above is equal to 
$$
\sum_{d=0}^\infty \sum_{\substack{\bs\in\N_0^d \\ |(\bs;0)|=s}} 
\lp J, w_{(\bkappa+\bs;0)} \rp
(-1)^s
\binom{k_1+s_1}{k_1}\cdots \binom{k_d+s_d}{k_d}.
$$
(Note: the prescribed condition $\bs\in\Z^d$ has been replaced 
with $\bs\in\N_0^d$ because of the a posteriori survivals of binomial factors.)
From this and \eqref{J-expansion} together with
$\lp J, w_{(\bkappa+\bs;0)} \rp=c_{(k_1+s_1,\dots,k_d+s_d;0)}$, 
we conclude the assertion.
\end{proof}

\begin{Note}[{\it Alternative proof of Lemma} \ref{lem2}] \label{Quest4.3}
In the right-hand side of Lemma \ref{lem2}, the quantity
$\binom{k_1+s_1}{k_1}\cdots \binom{k_d+s_d}{k_d}$ appearing there can also be interpreted as
the pairing
$\lp w_{(k_1,\dots,k_d;0)}\sha X^s, w_{(k_1+s_1,\dots,k_d+s_d;0)} \rp$.
Therefore, the assertion of this lemma is equivalent to the identity
\begin{equation}
\sum_{\btau\in\N_0^d} 
\lp\SS^{(\btau;0)},w_{(\bkappa+\bs;0)} \rp \cdot \lp\MM^{(\btau;0)}, w_{(\bkappa;0)}\!\cdot\!\! X^s\rp
=(-1)^s 
\lp w_{(\bkappa;0)}\sha X^s, w_{(\bkappa+\bs;0)} \rp
\label{alterLemma}
\end{equation}
for $\bkappa=(k_1,\dots,k_d)\in\N_0^d$, $\bs=(s_1,\dots,s_d)\in\Z^d$ 
satisfying $s=s_1+\cdots+s_d\ge 0$ and $\bkappa+\bs\in\N_0^d$.
We now give an alternative proof for it using the Magnus/demi-shuffle duality.
First, by Corollary \ref{coro}, we have 
$w_{(\bkappa;0)}=\sum_\bbr\la \MM^{(\bbr)},w_{(\bkappa;0)}\ra \SS^{(\bbr)}$ 
and $w_{(\bkappa+\bs;0)}=\sum_\bt\la \SS^{(\bt)},w_{(\bkappa+\bs;0)}\ra \MM^{(\bt)}$
so that the right-hand side of \eqref{alterLemma} can be written as
\begin{align}
\label{RHS4.2}
&(-1)^s 
\lp w_{(\bkappa;0)}\sha X^s, w_{(\bkappa+\bs;0)} \rp \\
&=(-1)^s \sum_{\bbr,\bt\in\N_0^{(\infty)}}
\lp \SS^{(\bbr)}\sha X^s, \MM^{(\bt)} \rp
\lp \MM^{(\bbr)} ,w_{(\bkappa;0)} \rp
\lp \SS^{(\bt)} ,w_{(\bkappa+\bs;0)}\rp
\notag \\
&=(-1)^s \sum_{\brho\in\N_0^d} 
\lp \MM^{(\brho;0)} ,w_{(\bkappa;0)}\rp
\lp  \SS^{(\brho;s)}, w_{(\bkappa+\bs;0)} \rp.
\notag
\end{align}
Here in the second equality, we use the fact that
$\lp  \MM^{(\bbr)}, w_{(\bkappa;0)} \rp$ 
survives only if $\bbr=(\brho;0)\in\N_0^{(\infty)}$ for some $\brho\in\N_0^d$
and then apply the duality (Theorem \ref{orthogonality})
to $\lp \SS^{(\bbr)}\sha X^s, \MM^{(\bt)} \rp$
with $\SS^{(\brho;0)}\sha X^s=\SS^{(\brho;s)}$ 
(cf. Definitions \ref{defMagnus} and \ref{demi-shuffle}).

On the other hand, in the left-hand side of  \eqref{alterLemma},
one observes that nontrivial terms of the summation arise
only from those $\btau=(\tau_1,\dots,\tau_d)\in\N_0^d$ with
$\sum_{i=1}^d \tau_i=s+\sum_{i=1}^d \kappa_i$ (a constant). 
But the last binomial factor in \eqref{MagnusBino} for
$ \lp\MM^{(\btau;0)}, w_{(\bkappa;0)}\!\cdot\!\! X^s\rp
=\magnus{(\tau_1,\dots,\tau_d;0)}{(\kappa_1,\dots,\kappa_d;s)}
$
equals $\binom{\tau_d}{s}$, which is non-zero only if
$\tau_d\ge s$.
Therefore, 
the summation~$\sum_\btau$ may be replaced by 
$\sum_\brho$ with $\brho=\btau-(\mathbf{0},s)$
in $\N_0^d$ (where $\mathbf{0}\in\N_0^{d-1}$ is the zero vector).
Thus,
the left-hand side of \eqref{alterLemma} can be written as
\begin{align}
\label{LHS4.2}
&\sum_{\btau\in\N_0^d} 
\lp\SS^{(\btau;0)},w_{(\bkappa+\bs;0)} \rp \cdot 
\lp\MM^{(\btau;0)}, w_{(\bkappa;0)}\!\cdot\!\! X^s\rp \\
&=
\sum_{\brho\in\N_0^d} 
\lp\SS^{(\brho+(\mathbf{0},s);0)},w_{(\bkappa+\bs;0)} \rp \cdot 
\lp\MM^{(\brho+(\mathbf{0},s);0)}, w_{(\bkappa;s)} \rp.
\notag
\end{align}
Comparing summands of the above \eqref{RHS4.2} and \eqref{LHS4.2} for 
individual $\brho\in\N_0^d$ in view of coefficients of 
monomial expansions of demi-shuffle/Magnus polynomials 
(Lemma~\ref{dS-mono} and~\eqref{Mag-Mono}), we reduce 
the formula \eqref{alterLemma} to the following elementary identity
for $\bkappa=(k_i)$, $\brho=(r_i)\in\N_0^d$ and $\bs=(s_i)\in\Z^d$
satisfying
$\sum_{i=1}^d k_i=\sum_{i=1}^d r_i$, $\bs+\bkappa\in \N_0^d$
and $s:=\sum_{i=1}^d s_i\ge 0$:
\begin{equation}
\binom{(\bkappa+\bs;0)}{(\brho;s)}
\magnus{(\brho;0)}{(\bkappa;0)}
=
(-1)^s
\binom{(\bkappa+\bs;0)}{(\brho+(\mathbf{0},s);0)}
\magnus{(\brho+(\mathbf{0},s);0)}{(\bkappa,s)}.
\end{equation} 
This is an immediate consequence of definitions
of these symbols $\{^*_* \}$, $(^*_*)$.
(Observe that only difference between the corresponding symbols
occurs from the last binomial coefficient in  \eqref{ArrayBino}
and \eqref{MagnusBino}.)
\qed
\end{Note}
% END AUTHOR BLOCK614--911

\begin{thebibliography}{99}
\bibitem{Ch21}Chapoton, Frédéric Zinbiel algebras and multiple zeta values, Doc. Math., Volume 27 (2022), pp. 519-533
\bibitem{Do10}Dokas, Ioannis Zinbiel algebras and commutative algebras with divided powers, Glasg. Math. J., Volume 52 (2010) no. 2, pp. 303-313
\bibitem{FP13}Foissy, Loïc; Patras, Frédéric Natural endomorphisms of shuffle algebras, Internat. J. Algebra Comput., Volume 23 (2013) no. 4, pp. 989-1009
\bibitem{F}Furusho, Hidekazu p-adic multiple zeta values. I. p-adic multiple polylogarithms and the p-adic KZ equation, Invent. Math., Volume 155 (2004) no. 2, pp. 253-286
\bibitem{GKP}Graham, Ronald L.; Knuth, Donald E.; Patashnik, Oren Concrete mathematics, Addison-Wesley Publishing Company, Reading, MA, 1994, xiv+657 pages (A foundation for computer science)
\bibitem{Minh2019}Hoang Ngoc Minh, Vincel On the solutions of the universal differential equation with three regular singularities (on solutions of KZ 3 ), Confluentes Math., Volume 11 (2019) no. 2, pp. 25-64
\bibitem{LM}Le, Thang Tu Quoc; Murakami, Jun Kontsevich’s integral for the Kauffman polynomial, Nagoya Math. J., Volume 142 (1996), pp. 39-65
\bibitem{Lo95}Loday, Jean-Louis Cup-product for Leibniz cohomology and dual Leibniz algebras, Math. Scand., Volume 77 (1995) no. 2, pp. 189-196
\bibitem{M37}Magnus, Wilhelm Über Beziehungen zwischen höheren Kommutatoren, J. Reine Angew. Math., Volume 177 (1937), pp. 105-115
\bibitem{MKS}Magnus, Wilhelm; Karrass, Abraham; Solitar, Donald Combinatorial group theory, Dover Publications, Inc., Mineola, NY, 2004, xii+444 pages (Presentations of groups in terms of generators and relations)
\bibitem{Maple}Maplesoft, a division of Waterloo Maple Inc.. Maple (2021 version) https://hadoop.apache.org
\bibitem{MPH}Minh, Hoang Ngoc; Petitot, Michel; Van Der Hoeven, Joris Shuffle algebra and polylogarithms, Discrete Math., Volume 225 (2000) no. 1-3, pp. 217-230 Formal power series and algebraic combinatorics (Toronto, ON, 1998)
\bibitem{N21}Nakamura, Hiroaki Some aspects of arithmetic functions in Grothendieck-Teichmüller theory, Oberwolfach Rep., Volume 18 (2021) no. 1, pp. 700-702 part of “Homotopic and geometric Galois theory. Abstracts from the workshop held March 7–13, 2021 (online meeting)”
\bibitem{Ree}Ree, Rimhak Lie elements and an algebra associated with shuffles, Ann. of Math. (2), Volume 68 (1958), pp. 210-220
\bibitem{Reu}Reutenauer, Christophe Free Lie algebras, London Mathematical Society Monographs. New Series, 7, The Clarendon Press, Oxford University Press, New York, 1993, xviii+269 pages (Oxford Science Publications)
\bibitem{Sch}Schützenberger, Marcel-Paul Sur une propriété combinatoire des demi-groupes libres, C. R. Acad. Sci. Paris, Volume 245 (1957), pp. 16-18
\end{thebibliography}
\end{document}
